\section*{Chapter \ref{ch:g11:the-eye} --- The Eye and Its Photoreceptors}

\begin{solution}{exo:g11:the-eye:1}
Cornea, pupil (through the iris), lens, vitreous body, then the
transparent layers of the retina (ganglion and bipolar \omterm{def:g10:cells-common-unit:cell}{cells}) to the
\omterm{def:g11:the-eye:photoreceptors}{photoreceptors} at the back.
\end{solution}

\begin{solution}{exo:g11:the-eye:2}
\omterm{def:g11:the-eye:photoreceptors}{Rods}: about 120 million, very sensitive, one pigment, night vision in
grey. \omterm{def:g11:the-eye:photoreceptors}{Cones}: about 6 million, need bright light, three pigments,
daylight and colour vision, sharpest at the fovea.
\end{solution}

\begin{solution}{exo:g11:the-eye:3}
The fovea is the central pit of the retina, made only of tightly
packed \omterm{def:g11:the-eye:photoreceptors}{cones}, where vision is sharpest. The blind spot is where the
optic nerve leaves, with no \omterm{def:g11:the-eye:photoreceptors}{photoreceptors}.
\end{solution}

\begin{solution}{exo:g11:the-eye:4}
S \omterm{def:g11:the-eye:photoreceptors}{cones} about 45\%, \omterm{def:g11:the-eye:photoreceptors}{rods} about 60\%, M \omterm{def:g11:the-eye:photoreceptors}{cones} about 20\%, L \omterm{def:g11:the-eye:photoreceptors}{cones} about
10\%.
\end{solution}

\begin{solution}{exo:g11:the-eye:5}
On the X \omterm{def:g11:cell-cycle-mitosis:chromosome}{chromosome}; the condition is X-linked \omterm{def:g11:genes-and-disease:genetic}{recessive}: expressed in
every male carrying the \omterm{def:g10:universal-dna:gene}{allele}, in females only with two copies, so
it is far commoner in men.
\end{solution}

\begin{solution}{exo:g11:the-eye:6}
At night only the \omterm{def:g11:the-eye:photoreceptors}{rods} respond, and they have a single pigment: they
report brightness, not wavelength, so no colour. The fovea has only
\omterm{def:g11:the-eye:photoreceptors}{cones}, too insensitive for the star; looked at sideways, the star's
image falls on \omterm{def:g11:the-eye:photoreceptors}{rods}, which detect it.
\end{solution}

\begin{solution}{exo:g11:the-eye:7}
Same colour: the ratios S:M:L are 0:0.6:1 in both. Different
brightness: light 2 excites every cone half as much. Colour is the
ratio; brightness is the total.
\end{solution}

\begin{solution}{exo:g11:the-eye:8}
Sons receive their X from the mother: all normal. Daughters receive
the father's X: all carriers, normal vision. Each carrier daughter's
sons are colour-blind with probability one half.
\end{solution}

\begin{solution}{exo:g11:the-eye:9}
A man needs one mutant X, a woman two. If the \omterm{def:g10:universal-dna:gene}{allele}'s frequency is
0.08, a woman carries two with probability $0.08^2 = 0.0064$, about
0.6\% --- sixteen times less than 8\%.
\end{solution}

\begin{solution}{exo:g11:the-eye:10}
Differences accumulate with time since a duplication: 4\% between L
and M means a recent copy; 57\% between either and S means a much
older one. The L/M duplication is the latest event, the S split far
earlier.
\end{solution}

\begin{solution}{exo:g11:the-eye:11}
The ancestor of Old World primates duplicated its single X-linked
\omterm{def:g11:the-eye:photoreceptors}{opsin} \omterm{def:g10:universal-dna:gene}{gene}; one copy shifted its peak to longer wavelengths, giving a
third cone type and a new dimension of colour. It was kept, probably
because it helped find ripe fruit and young leaves against foliage.
\end{solution}

\begin{solution}{exo:g11:the-eye:12}
Blues and yellows (and blue-green from green). The S \omterm{def:g10:universal-dna:gene}{gene} is on an
ordinary \omterm{def:g11:cell-cycle-mitosis:chromosome}{chromosome}, so both \omterm{def:g10:universal-dna:gene}{alleles} must be mutant in either sex ---
rare and sex-independent.
\end{solution}

\begin{solution}{exo:g11:the-eye:13}
$\num{2.5e-6}/\num{17e-3} \approx \num{1.5e-4}\,\unit{rad}$, about
half a minute of arc. At \qty{10}{m}: $\num{1.5e-4} \times 10 =
\qty{1.5}{mm}$.
\end{solution}

\begin{solution}{exo:g11:the-eye:14}
A male has one X, hence one \omterm{def:g10:universal-dna:gene}{allele}, hence one M/L-type cone: two cone
types with S. A female with two different \omterm{def:g10:universal-dna:gene}{alleles} expresses one in
some \omterm{def:g11:the-eye:photoreceptors}{cones} and the other in others: three types, hence three-colour
vision; a female with two identical \omterm{def:g10:universal-dna:gene}{alleles} has two types like a male.
\end{solution}

\begin{solution}{exo:g11:the-eye:15}
The ganglion \omterm{def:g10:cells-common-unit:cell}{cells} and the optic nerve are intact, so electrical
pulses delivered to them reach the visual brain as signals. Quality is
limited by the number of electrodes (hundreds, against a million
fibres and millions of receptors) and by the loss of the retina's own
sorting of the image.
\end{solution}

\begin{solution}{pb:g11:the-eye:1}
\textbf{1.} \qty{470}{nm}: S 40, M 35, L 18. \qty{530}{nm}: S 0, M 100,
L 80. \qty{590}{nm}: S 0, M 40, L 90. \qty{650}{nm}: S 0, M 2, L 12.

\textbf{2.} \qty{470}{nm} excites all three; \qty{650}{nm} essentially
only L.

\textbf{3.} With S and M only: \qty{590}{nm} gives (0, 40) and
\qty{650}{nm} gives (0, 2); \qty{530}{nm} gives (0, 100). The
\qty{590}{nm} (orange) and \qty{530}{nm} (green) lights differ only in
M intensity, like a dim and a bright green, and are the hard pair;
\qty{650}{nm} looks like a very dim light.

\textbf{4.} With S and L only: \qty{530}{nm} gives (0, 80) and
\qty{590}{nm} gives (0, 90) --- nearly identical: green and orange are
confused.

\textbf{5.} Both still detect red light, with the cone type they
keep; what is lost is the comparison between M and L responses, which
is what separates red from green.

\textbf{6.} $X^NX^c$: her father gave her his only X, which carries
the mutant \omterm{def:g10:universal-dna:gene}{allele}; her normal vision shows the other X is normal.

\textbf{7.} Sons: $X^NY$ normal or $X^cY$ colour-blind, one half each.
Daughters: $X^NX^N$ or $X^NX^c$, both normal vision, one half each
(the second carriers).

\textbf{8.} Ines is a carrier with probability $1/2$. A daughter is
colour-blind if she receives $X^c$ from both parents: from the father
certainly, from Ines with probability $1/2 \times 1/2$: $1/4$.

\textbf{9.} A colour-blind daughter has two mutant X \omterm{def:g10:universal-dna:information}{chromosomes}, one
from each parent; a father passes his only X to every daughter, so
his X is mutant and he is colour-blind.

\textbf{10.} Yes: her son received his X from her, and it carries the
mutant \omterm{def:g10:universal-dna:gene}{allele}; with normal vision, she was $X^NX^c$.

\textbf{11.} Rhodopsin branches off first; then S separates from the
L/M line; L and M split last: (rhodopsin, (S, (L, M))).

\textbf{12.} 4\% difference for 35 million years; 57\% would take
about $35 \times 57/4 \approx 500$ million years --- consistent with
the early \omterm{def:g10:common-ancestry:bodyplan}{vertebrate} origin (the relation is not exactly proportional
at large differences, so this is an order of magnitude).

\textbf{13.} The two nearly identical adjacent \omterm{def:g10:universal-dna:gene}{genes} can misalign
during the pairing of the two X \omterm{def:g10:universal-dna:information}{chromosomes}; an exchange then gives
one X with three copies and one with a single copy. A male receiving
the single-copy X has only one M/L-type cone: red--green colour
blindness.

\textbf{14.} Pairing relies on sequence similarity: two adjacent 96\%
identical \omterm{def:g10:universal-dna:gene}{genes} can pair out of register, \omterm{def:g10:universal-dna:gene}{gene} 1 with \omterm{def:g10:universal-dna:gene}{gene} 2, and
exchange unequally.

\textbf{15.} It occurred once, in the common ancestor of the Old World
primates, after their separation from other mammals and New World
monkeys, some 35 million years ago, and was inherited by all their
descendants.

\textbf{16.} $\num{2.5e-6}/\num{17e-3} \approx \num{1.5e-4}$ rad, about
$0.5'$.

\textbf{17.} $\num{1.5e-4} \times 6 \approx \qty{0.9}{mm}$: finer than
the \qty{1.7}{mm} strokes, which is why normal vision reads that line
comfortably (in practice optical blur brings the limit close to
$1'$).

\textbf{18.} Ten times coarser at least (100 \omterm{def:g11:the-eye:photoreceptors}{rods} share one signal,
about 10 by 10): several minutes of arc; in exchange, a hundred
receptors pool their light and detect far dimmer sources.

\textbf{19.} Twice the density is $\sqrt{2}$ closer spacing, and a
larger \omterm{def:g11:the-eye:parts}{eye} adds a longer focal length: about twice as fine.

\textbf{20.} S, M and L \omterm{def:g11:the-eye:photoreceptors}{cones}, each nearly blind to the far end of the
spectrum away from its peak; one tandem duplication of the X-linked
\omterm{def:g11:the-eye:photoreceptors}{opsin} \omterm{def:g10:universal-dna:gene}{gene}, 35 million years ago, gave primates the third; the fovea
resolves about half a minute of arc.
\end{solution}
